% ============================================================
% !TeX root = ../main.tex
%  Chapter 5 — Conclusions
% ============================================================
\chapter{Conclusions}
\label{ch:conclusions}

\noindent
This chapter closes the thesis. It summarises what was built and why, answers
the three research questions of \cref{sec:research-questions} on the evidence
of \cref{ch:results}, states the contributions, delimits the claims, and sets
out the work that should follow.

\section{Summary of the Work}
\label{sec:conclusion-summary}
\noindent
The thesis started from a geometric limitation of short-axis (SAX) cardiac MRI.
Each slice observes the left ventricle (LV) finely, at an in-plane pixel size
below $2$~mm, but the slices lie several millimetres apart, so the
three-dimensional wall on which clinical measurement depends is never observed
directly and has to be inferred. We treated that inference as the learning of a
continuous signed-distance representation. A phase-conditioned implicit model
was pre-trained on synthetic shapes sampled from a statistical shape model
(SSM) and fine-tuned on end-diastolic (ED) and end-systolic (ES) surfaces fitted
to ACDC and M\&Ms-2 segmentations. The model reads the sparse contour rings
through a global PointNet pathway and a local feature grid, predicts the
endocardial signed-distance field, and defines the epicardial field as that
field minus a strictly positive, directly supervised offset, so the two surfaces
cannot cross. Its output was compared with a radial-basis-function (RBF)
implicit fit and an SSM fit of the same rings, and myocardial wall thickness was
measured on all three geometries with four local estimators, localised on the
AHA 17-segment model, and contrasted between normal and hypertrophic patients.

\section{Answers to the Research Questions}
\label{sec:conclusion-research-questions}
\noindent
\textbf{RQ1} asked to what extent a continuous, watertight surface can be
reconstructed from sparse contours, and how such a reconstruction compares with
conventional implicit and shape-model alternatives. On the held-out ACDC and
M\&Ms-2 test patients, the model reproduces its input rings to within
$1.50$~mm on the endocardium and $1.96$~mm on the epicardium, which is the order
of the in-plane pixel size, and every surface it produces is watertight after
the shared repair stage. Its geometry lies closer to the RBF interpolant
(Chamfer distance $1.57$ and $1.90$~mm) than to the SSM fit ($2.08$ and
$2.59$~mm), and it responds to pathology more evenly than the population prior:
between the normal and hypertrophic subsets its mean contour distance moves from
$1.45$ to $1.60$~mm, whereas that of the SSM fit moves from $1.42$ to
$2.63$~mm. The comparison also shows where the three priors diverge. They agree
on the cavity, with a Dice coefficient of $0.88$ to $0.91$ and enclosed volumes
of $119$, $119$ and $127$~ml, but not on the wall, where the Dice coefficient
falls to $0.69$ to $0.77$ and the myocardial volume ranges from $88$ to
$136$~ml. Removing slices down to three per stack degrades all three methods by
a similar amount without breaking any of them. A closed, contour-faithful
surface can therefore be recovered from sparse SAX rings, but the wall
completed between the slices remains a property of the prior rather than of the
data, and the spread between methods is the honest measure of that uncertainty.

\textbf{RQ2} asked whether a single model can represent multiple physiological
states while preserving a physically plausible separation between the inner
and outer surfaces. One phase-conditioned network reconstructs both ED and ES
geometry, and because the epicardial field is the endocardial field minus a
positive offset, the ordering of the two surfaces is fixed by the architecture
rather than checked after the fact. On the reconstructed surfaces, the
Laplace-derived wall thickness rises from $9.21$~mm at ED to $13.87$~mm at ES
over the ACDC pathology subset, and all seventeen AHA segments thicken, with
the smallest relative change at the base and the largest in the mid-cavity and
apical segments. The evaluation did not, however, measure a pointwise minimum
separation on the extracted meshes, test for intersections after repair, or
report the ED--ES change in cavity volume, and the public SSM offers no ES
modes against which the systolic geometry could be compared. RQ2 therefore
receives a qualified answer: the model represents both phases with a wall that
thickens consistently between them, but positive separation is established by
construction of the fields and by regional means, not verified at every point
of the output.

\textbf{RQ3} asked how closely measurements derived from the reconstructed
geometry agree with those from reference clinical geometry, and how that
agreement depends on the measurement method. Both factors matter, and they act
in different ways. With the estimator held fixed, the reconstruction prior
shifts the absolute scale of the wall: patient-level mean thickness on the
model correlates with that on the RBF fit at $r = 0.99$ for the three non-ray
estimators, yet is $0.82$ to $1.16$~mm thicker, and against the SSM fit the
correlation falls to $0.72$ to $0.77$ with a bias of $2.53$ to $2.91$~mm. With
the geometry held fixed, the Yezzi--Prince method and the EDT boundary sum
agree with the Laplace field to within a bias below $1$~mm and an intraclass
correlation of $0.77$, whereas the SDF cone-ray estimator is biased by
$+3.31$~mm with limits of agreement of $\pm 12.74$~mm and an intraclass
correlation of $0.28$, because a straight ray leaves a curved wall where a
streamline does not. Strong correlation between geometries thus does not imply
agreement in millimetres; the reconstruction prior and the measurement path are
separate sources of offset of comparable size, and path-following estimators
are the appropriate choice on a curved wall. Because the RBF fit is a
smooth-surface comparator rather than an anatomical reference, these figures
bound the reproducibility of the measurement, not its accuracy.

The pathology-stratified analysis adds that the reconstructed geometry retains
patient-specific variation instead of collapsing towards an average shape. The
hypertrophic group is thicker in every AHA segment; its patient-level mean ED
thickness exceeds that of the normal group by $3.67$~mm (95\% CI
$[2.54,\ 4.67]$) and its thickest segment by $5.37$~mm ($[4.11,\ 6.56]$), both
with large effect sizes, whereas relative systolic thickening does not separate
the groups. These are descriptive separations within the model output on ten
hypertrophic patients, not a diagnostic result.

\section{Contributions}
\label{sec:conclusion-contributions}
\noindent
The thesis makes three contributions. The first is a data pipeline that brings
synthetic SSM meshes and clinical ED and ES segmentations from two public
datasets into one sample format, with contour rings, surface samples, dense
signed-distance supervision and wall-offset targets, so that a model can be
pre-trained on population geometry and fine-tuned on patient data without a
change of representation. The second is the phase-conditioned implicit model,
whose encoder pools endocardial and epicardial evidence separately before
merging it and complements the global code with a local feature grid read at
every query point, and whose decoder couples the epicardial field to the
endocardial field through a directly supervised positive offset. The third is
an evaluation protocol for sparse-contour reconstruction in the absence of 3D
ground truth: fidelity to the observed rings, symmetric agreement between
differently motivated reconstructions, self-consistency under slice
decimation, and a wall-thickness comparison that separates the effect of the
geometry from the effect of the estimator before the values are localised on
the AHA model and contrasted between patient groups. The protocol does not
depend on this particular network, and its central finding, that reconstruction
priors agree about the cavity and disagree about the wall, concerns the problem
rather than the model.

\section{Limitations}
\label{sec:conclusion-limitations}
\noindent
The answers above hold within the limits set out in \cref{sec:threats}. No
independent three-dimensional reference exists for these scans, so every
geometric figure in the thesis measures fidelity to the rings or agreement
between methods, never distance from the true anatomy, and every thickness
figure describes reproducibility rather than validated measurement. Because the
RBF comparator belongs to the same family as the surrogate surfaces used for
fine-tuning, its close agreement with the model is partly expected. Twenty
normal and ten hypertrophic ACDC patients support a directional observation
about pathology and no more, and the seventeen segment values of each patient
are repeated measures within one heart rather than independent samples. The SSM
comparison is confined to end-diastole, no ablation isolates the contribution
of synthetic pre-training, local conditioning, phase conditioning or the
positive-offset coupling, and the pre-training loss history was not retained.
What the evaluation offers is a careful characterisation of a working system,
not a validation of it.

\section{Future Work}
\label{sec:conclusion-future-work}
\noindent
The limitations fix the order of the next steps. The first need is an
anatomical reference. Higher-resolution volumetric acquisitions, or expert
three-dimensional surface and wall-thickness annotations on a larger external
cohort, would turn the agreement figures reported here into accuracy figures
and would show which of the three priors lies closest to the true wall. The
second is a wider and fully matched comparison, adding spline lofting, Poisson
reconstruction and learned alternatives under the same patients, physical
spacing, repair procedure and metrics, with pointwise error maps that localise
disagreement at the apex, the base and the most curved parts of the wall. The
third is the controlled ablation that this thesis did not run: removing
synthetic pre-training, local conditioning, phase conditioning and the
positive-offset coupling in turn would show which components the behaviour
depends on, and adding pointwise separation and intersection checks on the
extracted meshes would complete the answer to RQ2. Replacing the binary phase
label with a continuous position in the cardiac cycle would then allow the
representation to be tested for coherence in time as well as continuity in
space.

In sum, this thesis shows that a phase-conditioned implicit representation can
turn a sparse SAX stack into closed, paired LV surfaces that honour the
observed contours, thicken consistently between diastole and systole, and
support locally resolved wall-thickness analysis. It also shows, through the
disagreement between equally plausible reconstructions of the same rings, how
much of the myocardial wall is still inferred rather than observed, and
therefore how carefully any thickness derived from it must be read.
